\nonstopmode{}
\documentclass[a4paper]{book}
\usepackage[times,inconsolata,hyper]{Rd}
\usepackage{makeidx}
\makeatletter\@ifl@t@r\fmtversion{2018/04/01}{}{\usepackage[utf8]{inputenc}}\makeatother
% \usepackage{graphicx} % @USE GRAPHICX@
\makeindex{}
\begin{document}
\chapter*{}
\begin{center}
{\textbf{\huge Package `DDESONN'}}
\par\bigskip{\large \today}
\end{center}
\ifthenelse{\boolean{Rd@use@hyper}}{\hypersetup{pdftitle = {DDESONN: A Deep Dynamic Experimental Self-Organizing Neural Network Framework}}}{}
\ifthenelse{\boolean{Rd@use@hyper}}{\hypersetup{pdfauthor = {Mathew William Armitage Fok}}}{}
\begin{description}
\raggedright{}
\item[Type]\AsIs{Package}
\item[Title]\AsIs{A Deep Dynamic Experimental Self-Organizing Neural Network
Framework}
\item[Version]\AsIs{0.1.0}
\item[Description]\AsIs{High-level helpers for constructing, training, and evaluating Deep Dynamic Experimental Self-Organizing Neural Networks (DDESONN). The package wraps the legacy R6 implementations and exposes a modern API for modelling workflows. Evaluation utilities include ROC and Precision–Recall curves with default AUC/AUPRC annotation, configurable via ddesonn_run() and ddesonn_fit().}
\item[License]\AsIs{MIT + file LICENSE}
\item[Encoding]\AsIs{UTF-8}
\item[Roxygen]\AsIs{list(markdown = TRUE)}
\item[RoxygenNote]\AsIs{7.3.3}
\item[Depends]\AsIs{R (>= 4.1.0)}
\item[Imports]\AsIs{R6, stats, utils, dplyr, openxlsx, tidyr, pROC, PRROC,
reshape2, digest, ggplot2}
\item[Suggests]\AsIs{testthat, knitr, rmarkdown, foreach, quantmod, randomForest,
reticulate, zoo, readxl}
\item[VignetteBuilder]\AsIs{knitr}
\item[NeedsCompilation]\AsIs{no}
\item[Author]\AsIs{Mathew William Armitage Fok [aut, cre]}
\item[Maintainer]\AsIs{Mathew William Armitage Fok }\email{quiksilver67213@yahoo.com}\AsIs{}
\end{description}
\Rdcontents{Contents}
\HeaderA{DDESONN}{DDESONN: Deep Dynamic Experimental Self-Organizing Neural Network Framework}{DDESONN}
\keyword{internal}{DDESONN}
%
\begin{Description}
High-level helpers for constructing, training, and evaluating Deep Dynamic Experimental Self-Organizing Neural Networks.
\end{Description}
\HeaderA{DDESONN}{DDESONN: Deep Dynamic Experimental Self-Organizing Neural Network Framework}{DDESONN}
\aliasA{DDESONN-package}{DDESONN}{DDESONN.Rdash.package}
\keyword{internal}{DDESONN}
%
\begin{Description}
High-level helpers for constructing, training, and evaluating
Deep Dynamic Experimental Self-Organizing Neural Networks.
\end{Description}
%
\begin{Author}
\strong{Maintainer}: Mathew William Armitage Fok \email{quiksilver67213@yahoo.com}

\end{Author}
\HeaderA{ddesonn\_activations}{Activation functions (DDESONN)}{ddesonn.Rul.activations}
\aliasA{arctangent}{ddesonn\_activations}{arctangent}
\aliasA{bent\_identity}{ddesonn\_activations}{bent.Rul.identity}
\aliasA{bent\_relu}{ddesonn\_activations}{bent.Rul.relu}
\aliasA{bent\_sigmoid}{ddesonn\_activations}{bent.Rul.sigmoid}
\aliasA{bent\_swish}{ddesonn\_activations}{bent.Rul.swish}
\aliasA{binary\_activation}{ddesonn\_activations}{binary.Rul.activation}
\aliasA{custom\_activation}{ddesonn\_activations}{custom.Rul.activation}
\aliasA{custom\_bent\_piecewise}{ddesonn\_activations}{custom.Rul.bent.Rul.piecewise}
\aliasA{custom\_binary\_activation}{ddesonn\_activations}{custom.Rul.binary.Rul.activation}
\aliasA{elu}{ddesonn\_activations}{elu}
\aliasA{gaussian}{ddesonn\_activations}{gaussian}
\aliasA{gelu}{ddesonn\_activations}{gelu}
\aliasA{hard\_sigmoid}{ddesonn\_activations}{hard.Rul.sigmoid}
\aliasA{identity}{ddesonn\_activations}{identity}
\aliasA{inverse\_linear\_unit}{ddesonn\_activations}{inverse.Rul.linear.Rul.unit}
\aliasA{isrlu}{ddesonn\_activations}{isrlu}
\aliasA{leaky\_bent}{ddesonn\_activations}{leaky.Rul.bent}
\aliasA{leaky\_relu}{ddesonn\_activations}{leaky.Rul.relu}
\aliasA{leaky\_selu}{ddesonn\_activations}{leaky.Rul.selu}
\aliasA{maxout}{ddesonn\_activations}{maxout}
\aliasA{mish}{ddesonn\_activations}{mish}
\aliasA{parametric\_bent\_relu}{ddesonn\_activations}{parametric.Rul.bent.Rul.relu}
\aliasA{prelu}{ddesonn\_activations}{prelu}
\aliasA{relu}{ddesonn\_activations}{relu}
\aliasA{selu}{ddesonn\_activations}{selu}
\aliasA{sigmoid}{ddesonn\_activations}{sigmoid}
\aliasA{sigmoid\_binary}{ddesonn\_activations}{sigmoid.Rul.binary}
\aliasA{sigmoid\_sharp}{ddesonn\_activations}{sigmoid.Rul.sharp}
\aliasA{sinusoid}{ddesonn\_activations}{sinusoid}
\aliasA{softmax}{ddesonn\_activations}{softmax}
\aliasA{softplus}{ddesonn\_activations}{softplus}
\aliasA{swish}{ddesonn\_activations}{swish}
\aliasA{tanh}{ddesonn\_activations}{tanh}
\aliasA{tanh\_relu\_hybrid}{ddesonn\_activations}{tanh.Rul.relu.Rul.hybrid}
%
\begin{Description}
A collection of activation functions used by DDESONN.
These functions operate on numeric vectors/matrices and preserve shape.
\end{Description}
%
\begin{Usage}
\begin{verbatim}
binary_activation(x)

custom_binary_activation(x, threshold = -1.08)

custom_activation(z)

bent_identity(x)

relu(x)

softplus(x)

leaky_relu(x, alpha = 0.01)

elu(x, alpha = 1)

tanh(x)

sigmoid(x)

hard_sigmoid(x)

swish(x)

sigmoid_binary(x)

gelu(x)

selu(x, lambda = 1.0507, alpha = 1.67326)

mish(x)

prelu(x, alpha = 0.01)

softmax(z)

maxout(x, w1 = 0.5, b1 = 1, w2 = -0.5, b2 = 0.5)

bent_relu(x)

bent_sigmoid(x)

arctangent(x)

sinusoid(x)

gaussian(x)

isrlu(x, alpha = 1)

bent_swish(x)

parametric_bent_relu(x, beta = 1)

leaky_bent(x, alpha = 0.01)

inverse_linear_unit(x)

tanh_relu_hybrid(x)

custom_bent_piecewise(x, threshold = 0.5)

sigmoid_sharp(x, temp = 5)

leaky_selu(x, alpha = 0.01, lambda = 1.0507)

identity(x)
\end{verbatim}
\end{Usage}
%
\begin{Arguments}
\begin{ldescription}
\item[\code{x}] Numeric vector/matrix.

\item[\code{threshold}] Threshold for piecewise.

\item[\code{z}] Numeric vector/matrix.

\item[\code{alpha}] Leak factor.

\item[\code{lambda}] SELU lambda.

\item[\code{w1}] Weight 1.

\item[\code{b1}] Bias 1.

\item[\code{w2}] Weight 2.

\item[\code{b2}] Bias 2.

\item[\code{beta}] Beta parameter.

\item[\code{temp}] Temperature/sharpness.
\end{ldescription}
\end{Arguments}
%
\begin{Details}
Many functions coerce inputs to matrix form and preserve dimensions.
Some functions are experimental and may not be suitable for training.
\end{Details}
%
\begin{Value}
A matrix of 0/1 values with the same dimensions as \code{x}.

A matrix of 0/1 values with the same dimensions as \code{x}.

A matrix of 0/1 values with the same dimensions as \code{z}.

Bent identity transform.

ReLU transform.

Softplus transform.

Leaky ReLU transform.

ELU transform.

tanh transform.

Sigmoid transform.

Hard sigmoid transform.

Swish transform.

0/1 matrix.

GELU transform.

SELU transform.

Mish transform.

PReLU transform.

Row-wise softmax probabilities.

Elementwise max of two affine transforms.

Bent ReLU transform.

Bent sigmoid transform.

atan transform.

sin transform.

exp(-x\textasciicircum{}2).

ISRLU transform.

Bent swish transform.

Parametric bent ReLU transform.

Leaky bent transform.

x/(1+abs(x)).

Hybrid transform.

Piecewise transform.

Sharpened sigmoid.

Leaky SELU transform.

Base identity function.
\end{Value}
\HeaderA{ddesonn\_activation\_defaults}{Default activation sequences for DDESONN helpers}{ddesonn.Rul.activation.Rul.defaults}
%
\begin{Description}
Compute sensible activation functions for hidden and output
layers based on the modelling mode and stage (training or prediction).
\end{Description}
%
\begin{Usage}
\begin{verbatim}
ddesonn_activation_defaults(
  mode = c("binary", "multiclass", "regression"),
  hidden_sizes = NULL,
  stage = c("train", "predict")
)
\end{verbatim}
\end{Usage}
%
\begin{Arguments}
\begin{ldescription}
\item[\code{mode}] Problem mode. One of \code{"binary"}, \code{"multiclass"}, or \code{"regression"}.

\item[\code{hidden\_sizes}] Integer vector describing the hidden layer widths.

\item[\code{stage}] Stage for which activations are required. Either \code{"train"} or \code{"predict"}.
\end{ldescription}
\end{Arguments}
%
\begin{Value}
A list of activation functions suitable for passing into the
underlying R6 classes.
\end{Value}
%
\begin{Examples}
\begin{ExampleCode}
ddesonn_activation_defaults("binary", hidden_sizes = c(32, 16))
ddesonn_activation_defaults("regression", hidden_sizes = 64, stage = "predict")

\end{ExampleCode}
\end{Examples}
\HeaderA{ddesonn\_activation\_derivatives}{Activation derivatives (DDESONN)}{ddesonn.Rul.activation.Rul.derivatives}
\aliasA{arctangent\_derivative}{ddesonn\_activation\_derivatives}{arctangent.Rul.derivative}
\aliasA{bent\_identity\_derivative}{ddesonn\_activation\_derivatives}{bent.Rul.identity.Rul.derivative}
\aliasA{bent\_relu\_derivative}{ddesonn\_activation\_derivatives}{bent.Rul.relu.Rul.derivative}
\aliasA{bent\_sigmoid\_derivative}{ddesonn\_activation\_derivatives}{bent.Rul.sigmoid.Rul.derivative}
\aliasA{bent\_swish\_derivative}{ddesonn\_activation\_derivatives}{bent.Rul.swish.Rul.derivative}
\aliasA{binary\_activation\_derivative}{ddesonn\_activation\_derivatives}{binary.Rul.activation.Rul.derivative}
\aliasA{custom\_activation\_derivative}{ddesonn\_activation\_derivatives}{custom.Rul.activation.Rul.derivative}
\aliasA{custom\_bent\_piecewise\_derivative}{ddesonn\_activation\_derivatives}{custom.Rul.bent.Rul.piecewise.Rul.derivative}
\aliasA{custom\_binary\_activation\_derivative}{ddesonn\_activation\_derivatives}{custom.Rul.binary.Rul.activation.Rul.derivative}
\aliasA{elu\_derivative}{ddesonn\_activation\_derivatives}{elu.Rul.derivative}
\aliasA{gaussian\_derivative}{ddesonn\_activation\_derivatives}{gaussian.Rul.derivative}
\aliasA{gelu\_derivative}{ddesonn\_activation\_derivatives}{gelu.Rul.derivative}
\aliasA{hard\_sigmoid\_derivative}{ddesonn\_activation\_derivatives}{hard.Rul.sigmoid.Rul.derivative}
\aliasA{identity\_derivative}{ddesonn\_activation\_derivatives}{identity.Rul.derivative}
\aliasA{inverse\_linear\_unit\_derivative}{ddesonn\_activation\_derivatives}{inverse.Rul.linear.Rul.unit.Rul.derivative}
\aliasA{isrlu\_derivative}{ddesonn\_activation\_derivatives}{isrlu.Rul.derivative}
\aliasA{leaky\_bent\_derivative}{ddesonn\_activation\_derivatives}{leaky.Rul.bent.Rul.derivative}
\aliasA{leaky\_relu\_derivative}{ddesonn\_activation\_derivatives}{leaky.Rul.relu.Rul.derivative}
\aliasA{leaky\_selu\_derivative}{ddesonn\_activation\_derivatives}{leaky.Rul.selu.Rul.derivative}
\aliasA{maxout\_derivative}{ddesonn\_activation\_derivatives}{maxout.Rul.derivative}
\aliasA{mish\_derivative}{ddesonn\_activation\_derivatives}{mish.Rul.derivative}
\aliasA{parametric\_bent\_relu\_derivative}{ddesonn\_activation\_derivatives}{parametric.Rul.bent.Rul.relu.Rul.derivative}
\aliasA{prelu\_derivative}{ddesonn\_activation\_derivatives}{prelu.Rul.derivative}
\aliasA{relu\_derivative}{ddesonn\_activation\_derivatives}{relu.Rul.derivative}
\aliasA{selu\_derivative}{ddesonn\_activation\_derivatives}{selu.Rul.derivative}
\aliasA{sigmoid\_binary\_derivative}{ddesonn\_activation\_derivatives}{sigmoid.Rul.binary.Rul.derivative}
\aliasA{sigmoid\_derivative}{ddesonn\_activation\_derivatives}{sigmoid.Rul.derivative}
\aliasA{sigmoid\_sharp\_derivative}{ddesonn\_activation\_derivatives}{sigmoid.Rul.sharp.Rul.derivative}
\aliasA{sinusoid\_derivative}{ddesonn\_activation\_derivatives}{sinusoid.Rul.derivative}
\aliasA{softmax\_derivative}{ddesonn\_activation\_derivatives}{softmax.Rul.derivative}
\aliasA{softplus\_derivative}{ddesonn\_activation\_derivatives}{softplus.Rul.derivative}
\aliasA{swish\_derivative}{ddesonn\_activation\_derivatives}{swish.Rul.derivative}
\aliasA{tanh\_derivative}{ddesonn\_activation\_derivatives}{tanh.Rul.derivative}
\aliasA{tanh\_relu\_hybrid\_derivative}{ddesonn\_activation\_derivatives}{tanh.Rul.relu.Rul.hybrid.Rul.derivative}
%
\begin{Description}
A collection of derivative functions corresponding to DDESONN activation functions.
\end{Description}
%
\begin{Usage}
\begin{verbatim}
binary_activation_derivative(x)

custom_binary_activation_derivative(x, threshold = -1.08)

custom_activation_derivative(z)

bent_identity_derivative(x)

relu_derivative(x)

softplus_derivative(x)

leaky_relu_derivative(x, alpha = 0.01)

elu_derivative(x, alpha = 1)

tanh_derivative(x)

sigmoid_derivative(x)

hard_sigmoid_derivative(x)

swish_derivative(x)

sigmoid_binary_derivative(x)

gelu_derivative(x)

selu_derivative(x, lambda = 1.0507, alpha = 1.67326)

mish_derivative(x)

maxout_derivative(x, w1 = 0.5, b1 = 1, w2 = -0.5, b2 = 0.5)

prelu_derivative(x, alpha = 0.01)

softmax_derivative(x)

bent_relu_derivative(x)

bent_sigmoid_derivative(x)

arctangent_derivative(x)

sinusoid_derivative(x)

gaussian_derivative(x)

isrlu_derivative(x, alpha = 1)

bent_swish_derivative(x)

parametric_bent_relu_derivative(x, beta = 1)

leaky_bent_derivative(x, alpha = 0.01)

inverse_linear_unit_derivative(x)

tanh_relu_hybrid_derivative(x)

custom_bent_piecewise_derivative(x, threshold = 0.5)

sigmoid_sharp_derivative(x, temp = 5)

leaky_selu_derivative(x, alpha = 0.01, lambda = 1.0507)

identity_derivative(x)
\end{verbatim}
\end{Usage}
%
\begin{Arguments}
\begin{ldescription}
\item[\code{x}] Numeric vector/matrix.

\item[\code{threshold}] Threshold.

\item[\code{z}] Numeric vector/matrix.

\item[\code{alpha}] Leak factor.

\item[\code{lambda}] SELU lambda.

\item[\code{w1}] Weight 1.

\item[\code{b1}] Bias 1.

\item[\code{w2}] Weight 2.

\item[\code{b2}] Bias 2.

\item[\code{beta}] Beta parameter.

\item[\code{temp}] Temperature/sharpness.
\end{ldescription}
\end{Arguments}
%
\begin{Value}
Vector of zeros.

Vector of zeros.

Zero array with same dims.

Derivative matrix.

Derivative matrix.

Derivative matrix.

Derivative matrix.

Derivative matrix.

Derivative matrix.

Derivative matrix.

Derivative matrix.

Derivative matrix.

Vector of zeros.

Derivative matrix.

Derivative matrix.

Derivative matrix.

Gradient of active affine branch.

Derivative matrix.

Derivative matrix (approx).

Derivative matrix.

Derivative matrix.

Derivative matrix.

Derivative matrix.

Derivative matrix.

Derivative matrix.

Derivative matrix.

Derivative matrix.

Derivative matrix.

Derivative matrix.

Derivative matrix.

Derivative matrix.

Derivative matrix.

Derivative matrix/vector.

Matrix of ones with same dimensions as \code{x}.
\end{Value}
\HeaderA{ddesonn\_artifacts\_root}{Resolve the writable artifacts root for DDESONN.}{ddesonn.Rul.artifacts.Rul.root}
%
\begin{Description}
Defaults to a user-writable data directory (via tools::R\_user\_dir()) so
runtime outputs never write into the installed package tree.
\end{Description}
%
\begin{Usage}
\begin{verbatim}
ddesonn_artifacts_root(output_root = NULL)
\end{verbatim}
\end{Usage}
%
\begin{Arguments}
\begin{ldescription}
\item[\code{output\_root}] Optional base directory for artifacts. When NULL, a
user-scoped directory is selected automatically.
\end{ldescription}
\end{Arguments}
%
\begin{Details}
Override order (first non-empty wins):
\begin{enumerate}

\item{} output\_root argument
\item{} Sys.getenv("DDESONN\_ARTIFACTS\_ROOT")
\item{} getOption("DDESONN\_OUTPUT\_ROOT")

\end{enumerate}

\end{Details}
%
\begin{Value}
Absolute path to the artifacts directory (created if missing).
\end{Value}
\HeaderA{ddesonn\_debug\_state}{Inspect internal DDESONN debug state}{ddesonn.Rul.debug.Rul.state}
%
\begin{Description}
Returns a named list snapshot of objects currently stored in the private package debug/state environment.
\end{Description}
%
\begin{Usage}
\begin{verbatim}
ddesonn_debug_state()
\end{verbatim}
\end{Usage}
%
\begin{Value}
A named list of objects from DDESONN internal state.
\end{Value}
\HeaderA{ddesonn\_dropout\_defaults}{Default dropout configuration}{ddesonn.Rul.dropout.Rul.defaults}
%
\begin{Description}
Produce a simple dropout configuration matching the supplied
hidden layer sizes.
\end{Description}
%
\begin{Usage}
\begin{verbatim}
ddesonn_dropout_defaults(hidden_sizes)
\end{verbatim}
\end{Usage}
%
\begin{Arguments}
\begin{ldescription}
\item[\code{hidden\_sizes}] Integer vector describing the hidden layer widths.
\end{ldescription}
\end{Arguments}
%
\begin{Value}
A list of dropout rates for each hidden layer.
\end{Value}
%
\begin{Examples}
\begin{ExampleCode}
ddesonn_dropout_defaults(c(64, 32))

\end{ExampleCode}
\end{Examples}
\HeaderA{ddesonn\_fit}{Fit a \code{ddesonn\_model} with tidy inputs}{ddesonn.Rul.fit}
%
\begin{Description}
Train a \code{ddesonn\_model} (backed by \code{DDESONN}) using matrices or
data frames, handling label coercion, validation data, and training control
defaults.
\end{Description}
%
\begin{Usage}
\begin{verbatim}
ddesonn_fit(
  model,
  x,
  y,
  validation = NULL,
  self_org = NULL,
  ...,
  verbose = FALSE,
  verboseLow = FALSE,
  debug = FALSE
)
\end{verbatim}
\end{Usage}
%
\begin{Arguments}
\begin{ldescription}
\item[\code{model}] A model created by \code{\LinkA{ddesonn\_model()}{ddesonn.Rul.model}}.

\item[\code{x}] Training features.

\item[\code{y}] Training targets/labels.

\item[\code{validation}] Optional list containing \code{x} and \code{y} elements for validation.

\item[\code{self\_org}] Optional logical override for the legacy self-organization
phase. \code{TRUE} enables \code{self\_organize()} during training and \code{FALSE}
disables it. \code{NULL} keeps the configured default (\code{self\_org = FALSE} in
\code{\LinkA{ddesonn\_training\_defaults()}{ddesonn.Rul.training.Rul.defaults}}). Self-organization acts on input-space
topology error (how well neighborhood structure is organized), not on the
supervised prediction-loss term.

\item[\code{...}] Named overrides for entries in \code{\LinkA{ddesonn\_training\_defaults()}{ddesonn.Rul.training.Rul.defaults}}.

\item[\code{verbose}] Logical; emit detailed progress output when TRUE.

\item[\code{verboseLow}] Logical; emit important progress output when TRUE.

\item[\code{debug}] Logical; emit debug diagnostics when TRUE.
\end{ldescription}
\end{Arguments}
%
\begin{Value}
The trained model (invisibly). The underlying R6 object is modified
in-place and the last training result is stored under \code{model\$last\_training}.
\end{Value}
%
\begin{SeeAlso}
\LinkA{DDESONN}{DDESONN}
\end{SeeAlso}
%
\begin{Examples}
\begin{ExampleCode}
data <- mtcars
x <- data[, c("disp", "hp", "wt", "qsec", "drat")]
y <- data$am
model <- ddesonn_model(input_size = ncol(x), output_size = 1, hidden_sizes = 8)
ddesonn_fit(model, x, y, num_epochs = 1, lr = 0.05, validation_metrics = FALSE)

# Regression example (mtcars) with explicit scheduler controls.
# If you do NOT want LR decay, set lr_decay_rate = 1.0.
reg_x <- mtcars[, c("disp", "hp", "wt", "qsec", "drat")]
reg_y <- mtcars$mpg
reg_model <- ddesonn_model(
  input_size = ncol(reg_x), # number of input features
  output_size = 1, # one numeric target
  hidden_sizes = c(16, 8), # hidden-layer widths
  classification_mode = "regression" # problem type
)
ddesonn_fit(
  model = reg_model, # model object from ddesonn_model()
  x = reg_x, # training predictors
  y = reg_y, # training target
  num_epochs = 10, # training epochs
  lr = 0.05, # initial learning rate
  lr_decay_rate = 0.5, # decay multiplier (use 1.0 to disable)
  lr_decay_epoch = 20L, # decay step interval in epochs
  lr_min = 1e-5, # lower bound for learning rate
  validation_metrics = FALSE # disable validation metric pass in this example
)

\end{ExampleCode}
\end{Examples}
\HeaderA{ddesonn\_model}{Create a high-level DDESONN model wrapper}{ddesonn.Rul.model}
%
\begin{Description}
Initialise a \code{ddesonn\_model} (R6) instance backed by the legacy
\code{DDESONN} class, while handling sensible defaults for activations and node
counts.
\end{Description}
%
\begin{Usage}
\begin{verbatim}
ddesonn_model(
  input_size,
  output_size,
  hidden_sizes = c(64, 32),
  num_networks = 1L,
  lambda = 0.00028,
  classification_mode = c("binary", "multiclass", "regression"),
  ML_NN = TRUE,
  activation_functions = NULL,
  activation_functions_predict = NULL,
  init_method = "he",
  custom_scale = 1,
  N = NULL,
  ensembles = NULL,
  ensemble_number = 0L,
  verbose = FALSE,
  verboseLow = FALSE,
  debug = FALSE
)
\end{verbatim}
\end{Usage}
%
\begin{Arguments}
\begin{ldescription}
\item[\code{input\_size}] Number of input features.

\item[\code{output\_size}] Number of outputs.

\item[\code{hidden\_sizes}] Integer vector describing hidden layer widths.

\item[\code{num\_networks}] Number of SONN members to initialise within the ensemble.

\item[\code{lambda}] Regularisation strength.

\item[\code{classification\_mode}] Problem mode: \code{"binary"}, \code{"multiclass"}, or \code{"regression"}.

\item[\code{ML\_NN}] Logical; whether to initialise a multi-layer SONN.

\item[\code{activation\_functions}] Optional list of activation functions for training.

\item[\code{activation\_functions\_predict}] Optional list of activation functions used during prediction.

\item[\code{init\_method}] Weight initialisation scheme passed to the legacy constructor.

\item[\code{custom\_scale}] Optional scaling factor for the initialiser.

\item[\code{N}] Optional total node count. If omitted it is inferred from the architecture.

\item[\code{ensembles}] Optional pre-existing ensemble container.

\item[\code{ensemble\_number}] Identifier used when combining multiple ensembles.

\item[\code{verbose}] Logical; emit detailed progress output when TRUE.

\item[\code{verboseLow}] Logical; emit important progress output when TRUE.

\item[\code{debug}] Logical; emit debug diagnostics when TRUE.
\end{ldescription}
\end{Arguments}
%
\begin{Value}
A \code{ddesonn\_model} (R6) instance ready for training.
\end{Value}
%
\begin{SeeAlso}
\LinkA{DDESONN}{DDESONN}
\end{SeeAlso}
%
\begin{Examples}
\begin{ExampleCode}
model <- ddesonn_model(
  input_size = 5,
  output_size = 1,
  hidden_sizes = c(32, 16),
  classification_mode = "binary"
)

\end{ExampleCode}
\end{Examples}
\HeaderA{ddesonn\_optimizer\_options}{Supported optimizer identifiers}{ddesonn.Rul.optimizer.Rul.options}
%
\begin{Description}
List the optimizer strings understood by the legacy DDESONN
training loop.
\end{Description}
%
\begin{Usage}
\begin{verbatim}
ddesonn_optimizer_options()
\end{verbatim}
\end{Usage}
%
\begin{Value}
A character vector of supported optimiser identifiers.
\end{Value}
%
\begin{Examples}
\begin{ExampleCode}
ddesonn_optimizer_options()

\end{ExampleCode}
\end{Examples}
\HeaderA{ddesonn\_plots\_dir}{Resolve the plots directory inside the artifacts root.}{ddesonn.Rul.plots.Rul.dir}
%
\begin{Description}
Resolve the plots directory inside the artifacts root.
\end{Description}
%
\begin{Usage}
\begin{verbatim}
ddesonn_plots_dir(output_root = NULL)
\end{verbatim}
\end{Usage}
%
\begin{Arguments}
\begin{ldescription}
\item[\code{output\_root}] Optional base directory for artifacts. When NULL, a
user-scoped directory is selected automatically.
\end{ldescription}
\end{Arguments}
%
\begin{Value}
Absolute path to the plots directory.
\end{Value}
\HeaderA{ddesonn\_predict}{Generate predictions from a fitted \code{ddesonn\_model}}{ddesonn.Rul.predict}
%
\begin{Description}
Internal prediction engine / forward-pass primitive that
produces ensemble or per-model predictions from a trained \code{ddesonn\_model}.
For user-facing inference, prefer \code{\LinkA{predict.ddesonn\_model()}{predict.ddesonn.Rul.model}}, which wraps
this helper to provide a stable API for \code{type}/\code{aggregate}/\code{threshold}
handling and return shapes.
Multiclass note: For multiclass classification, y should be encoded as integer class indices 1..K (or a one-hot matrix whose columns follow the model’s class order), otherwise accuracy comparisons may be incorrect.
\end{Description}
%
\begin{Usage}
\begin{verbatim}
ddesonn_predict(
  model,
  new_data,
  aggregate = c("mean", "median", "none"),
  type = c("response", "class"),
  threshold = NULL,
  verbose = FALSE,
  verboseLow = FALSE,
  debug = FALSE
)
\end{verbatim}
\end{Usage}
%
\begin{Arguments}
\begin{ldescription}
\item[\code{model}] A trained model produced by \code{\LinkA{ddesonn\_model()}{ddesonn.Rul.model}}.

\item[\code{new\_data}] New feature matrix or data frame.

\item[\code{aggregate}] Aggregation strategy across ensemble members. One of
\code{"mean"}, \code{"median"}, or \code{"none"}.

\item[\code{type}] Prediction type. \code{"response"} returns numeric predictions,
while \code{"class"} applies thresholding for classification problems.

\item[\code{threshold}] Optional threshold override when \code{type = "class"}.

\item[\code{verbose}] Logical; emit detailed progress output when TRUE.

\item[\code{verboseLow}] Logical; emit important progress output when TRUE.

\item[\code{debug}] Logical; emit debug diagnostics when TRUE.
\end{ldescription}
\end{Arguments}
%
\begin{Value}
A list containing the aggregated prediction matrix and the
per-model outputs when \code{aggregate = "none"}.
\end{Value}
%
\begin{SeeAlso}
\LinkA{DDESONN}{DDESONN}
\end{SeeAlso}
%
\begin{Examples}
\begin{ExampleCode}
data <- mtcars
x <- data[, c("disp", "hp", "wt", "qsec", "drat")]
y <- data$am
model <- ddesonn_model(input_size = ncol(x), output_size = 1, hidden_sizes = 8)
ddesonn_fit(model, x, y, num_epochs = 1, lr = 0.05, validation_metrics = FALSE)
preds <- ddesonn_predict(model, x)
head(preds$prediction)

\end{ExampleCode}
\end{Examples}
\HeaderA{ddesonn\_run}{Run DDESONN across common ensemble scenarios.}{ddesonn.Rul.run}
%
\begin{Description}
This helper re-creates the four orchestration modes that previously lived in
\code{TestDDESONN.R}:
\end{Description}
%
\begin{Usage}
\begin{verbatim}
ddesonn_run(
  x,
  y,
  classification_mode = c("binary", "multiclass", "regression"),
  hidden_sizes = c(64, 32),
  seeds = 1L,
  do_ensemble = FALSE,
  num_networks = if (isTRUE(do_ensemble)) 3L else 1L,
  num_temp_iterations = 0L,
  validation = NULL,
  x_valid = NULL,
  y_valid = NULL,
  model_overrides = list(),
  training_overrides = list(),
  temp_overrides = NULL,
  prediction_data = NULL,
  test = NULL,
  x_test = NULL,
  y_test = NULL,
  prediction_type = c("response", "class"),
  aggregate = c("mean", "median", "none"),
  seed_aggregate = c("mean", "median", "none"),
  threshold = NULL,
  output_root = NULL,
  plot_controls = NULL,
  save_models = TRUE,
  verbose = FALSE,
  verboseLow = FALSE,
  debug = FALSE
)
\end{verbatim}
\end{Usage}
%
\begin{Arguments}
\begin{ldescription}
\item[\code{x}] Training features (the training split) as a data frame, matrix, or tibble.

\item[\code{y}] Training labels/targets (the training split).

\item[\code{classification\_mode}] Overall problem mode. One of \code{"binary"},
\code{"multiclass"}, or \code{"regression"}.

\item[\code{hidden\_sizes}] Integer vector describing the hidden layer widths.

\item[\code{seeds}] Integer vector of seeds. A separate model (or ensemble stack) is
trained for each seed.

\item[\code{do\_ensemble}] Logical flag selecting the ensemble container modes
(scenarios C/D). When \code{FALSE}, scenarios A/B are executed.

\item[\code{num\_networks}] Number of ensemble members inside each
\code{\LinkA{ddesonn\_model()}{ddesonn.Rul.model}} instance.

\item[\code{num\_temp\_iterations}] Number of TEMP iterations to run when
\code{do\_ensemble = TRUE} (scenario D). Ignored otherwise.

\item[\code{validation}] Optional validation list with elements \code{x} and \code{y}.

\item[\code{x\_valid}] Optional validation features. Overrides \code{validation\$x} when set.

\item[\code{y\_valid}] Optional validation labels. Overrides \code{validation\$y} when set.

\item[\code{model\_overrides}] Named list forwarded to \code{\LinkA{ddesonn\_model()}{ddesonn.Rul.model}} allowing
custom architectures.

\item[\code{training\_overrides}] Named list forwarded to \code{\LinkA{ddesonn\_fit()}{ddesonn.Rul.fit}} for the
main run(s). Any argument accepted by \code{\LinkA{ddesonn\_fit()}{ddesonn.Rul.fit}} may be provided here.
Unspecified values fall back to \code{\LinkA{ddesonn\_training\_defaults()}{ddesonn.Rul.training.Rul.defaults}} for the given
\code{classification\_mode} and \code{hidden\_sizes}. See \strong{Details} and the example
showing how to inspect defaults.

\item[\code{temp\_overrides}] Optional named list forwarded to \code{\LinkA{ddesonn\_fit()}{ddesonn.Rul.fit}} for
TEMP iterations. Defaults to \code{training\_overrides}. Use this when TEMP
candidates should train differently than the main model.

\item[\code{prediction\_data}] Optional features for prediction. When supplied,
predictions are computed for each seed/iteration.

\item[\code{test}] Optional test list with elements \code{x} and \code{y}. When supplied,
the final model computes test metrics (loss and, for classification,
accuracy) and stores them in \code{result\$test\_metrics}. The run history
(\code{result\$history}) mirrors the training metadata (train/validation losses)
and appends \code{test\_loss} when test data is provided.

\item[\code{x\_test}] Optional test features. Overrides \code{test\$x} when set.

\item[\code{y\_test}] Optional test labels. Overrides \code{test\$y} when set.

\item[\code{prediction\_type}] Passed to \code{\LinkA{ddesonn\_predict()}{ddesonn.Rul.predict}}.

\item[\code{aggregate}] Aggregation strategy within a single model (across ensemble
members).

\item[\code{seed\_aggregate}] Aggregation strategy across seeds. Set to \code{"none"} to
keep per-seed prediction matrices.

\item[\code{threshold}] Optional threshold override for classification prediction.

\item[\code{output\_root}] Optional directory where legacy-style artifacts are
written. When \code{NULL} (default) no files are created.

\item[\code{plot\_controls}] Optional list passed through to \code{\LinkA{ddesonn\_fit()}{ddesonn.Rul.fit}} as
\code{plot\_controls}. Use this to enable/disable specific report plots or
diagnostics (for example, evaluation report settings). The supported
structure is defined by \code{\LinkA{ddesonn\_fit()}{ddesonn.Rul.fit}}; this function does not create
defaults.

\item[\code{save\_models}] Logical; if \code{TRUE} (default) individual models are
persisted when \code{output\_root} is supplied.

\item[\code{verbose}] Logical; emit detailed progress output when TRUE.

\item[\code{verboseLow}] Logical; emit important progress output when TRUE.

\item[\code{debug}] Logical; emit debug diagnostics when TRUE.
\end{ldescription}
\end{Arguments}
%
\begin{Details}
\begin{itemize}

\item{} Scenario A – single model (\code{do\_ensemble = FALSE}, \code{num\_networks = 1}).
\item{} Scenario B – single run with multiple members inside a single model
(\code{do\_ensemble = FALSE}, \code{num\_networks > 1}).
\item{} Scenario C – main ensemble container (\code{do\_ensemble = TRUE},
\code{num\_temp\_iterations = 0}).
\item{} Scenario D – main ensemble plus TEMP iterations (\code{do\_ensemble = TRUE},
\code{num\_temp\_iterations > 0}).

\end{itemize}


The function accepts a training set, optional validation data, and optional
prediction features. It repeatedly instantiates \code{\LinkA{ddesonn\_model()}{ddesonn.Rul.model}} objects,
fits them with \code{\LinkA{ddesonn\_fit()}{ddesonn.Rul.fit}}, and (when requested) calls
\code{\LinkA{ddesonn\_predict()}{ddesonn.Rul.predict}} to surface aggregated predictions.

\strong{Discovering available training overrides}

\code{training\_overrides} is a direct pass-through to \code{\LinkA{ddesonn\_fit()}{ddesonn.Rul.fit}}. To see the
baseline defaults used by \code{ddesonn\_run()}, call:

\code{ddesonn\_training\_defaults(classification\_mode, hidden\_sizes)}

To see all tunable training arguments, see \code{?ddesonn\_fit}.
\end{Details}
%
\begin{Value}
A list (classed as \code{"ddesonn\_run\_result"}) containing the
configuration, per-seed models, and optional prediction summaries.
\end{Value}
%
\begin{Examples}
\begin{ExampleCode}

# ============================================================
# DDESONN — FULL example using package data in inst/extdata
# (binary classification; train/valid/test split; scale train-only)
# ============================================================

library(DDESONN)

set.seed(111)

# ------------------------------------------------------------
# 1) Locate package extdata folder (robust across check/install)  #$$$$$$$$$$$$$
# ------------------------------------------------------------
ext_dir <- system.file("extdata", package = "DDESONN")
if (!nzchar(ext_dir)) {
  stop("Could not find DDESONN extdata folder. Is the package installed?",
       call. = FALSE)
}

# ------------------------------------------------------------
# 1b) Find CSVs (recursive + check-dir edge cases)               #$$$$$$$$$$$$$
# ------------------------------------------------------------
csvs <- list.files(
  ext_dir,
  pattern = "\\\\.csv$",
  full.names = TRUE,
  recursive = TRUE
)

# Defensive fallback for rare nested layouts
if (!length(csvs)) {                                             #$$$$$$$$$$$$$
  ext_dir2 <- file.path(ext_dir, "inst", "extdata")               #$$$$$$$$$$$$$
  if (dir.exists(ext_dir2)) {                                    #$$$$$$$$$$$$$
    csvs <- list.files(
      ext_dir2,
      pattern = "\\\\.csv$",
      full.names = TRUE,
      recursive = TRUE
    )
  }
}

if (!length(csvs)) {
  message(sprintf(
    "No .csv files found under: %s — skipping example.",
    ext_dir
  ))
} else {

  hf_path <- file.path(ext_dir, "heart_failure_clinical_records.csv")
  data_path <- if (file.exists(hf_path)) hf_path else csvs[[1]]

  cat("[extdata] using:", data_path, "\\n")

# ------------------------------------------------------------
# 2) Load data
# ------------------------------------------------------------
df <- read.csv(data_path)

# Prefer DEATH_EVENT if present; otherwise infer a binary target
target_col <- if ("DEATH_EVENT" %in% names(df)) {
  "DEATH_EVENT"
} else {
  cand <- names(df)[vapply(df, function(col) {
    v <- suppressWarnings(as.numeric(col))
    if (all(is.na(v))) return(FALSE)
    u <- unique(v[is.finite(v)])
    length(u) <= 2 && all(sort(u) %in% c(0, 1))
  }, logical(1))]
  if (!length(cand)) {
    stop(
      "Could not infer a binary target column. ",
      "Provide a binary CSV in extdata or rename target to DEATH_EVENT.",
      call. = FALSE
    )
  }
  cand[[1]]
}

cat("[data] target_col =", target_col, "\\n")

# ------------------------------------------------------------
# 3) Build X and y
# ------------------------------------------------------------
y_all <- matrix(as.integer(df[[target_col]]), ncol = 1)

x_df <- df[, setdiff(names(df), target_col), drop = FALSE]
x_all <- as.matrix(x_df)
storage.mode(x_all) <- "double"

# ------------------------------------------------------------
# 4) Split 70 / 15 / 15
# ------------------------------------------------------------
n <- nrow(x_all)
idx <- sample.int(n)

n_train <- floor(0.70 * n)
n_valid <- floor(0.15 * n)

i_tr <- idx[1:n_train]
i_va <- idx[(n_train + 1):(n_train + n_valid)]
i_te <- idx[(n_train + n_valid + 1):n]

x_train <- x_all[i_tr, , drop = FALSE]
y_train <- y_all[i_tr, , drop = FALSE]

x_valid <- x_all[i_va, , drop = FALSE]
y_valid <- y_all[i_va, , drop = FALSE]

x_test  <- x_all[i_te, , drop = FALSE]
y_test  <- y_all[i_te, , drop = FALSE]

cat(sprintf("[split] train=%d valid=%d test=%d\\n",
            nrow(x_train), nrow(x_valid), nrow(x_test)))

# ------------------------------------------------------------
# 5) Scale using TRAIN stats only (no leakage)
# ------------------------------------------------------------
x_train_s <- scale(x_train)
ctr <- attr(x_train_s, "scaled:center")
scl <- attr(x_train_s, "scaled:scale")
scl[!is.finite(scl) | scl == 0] <- 1

x_valid_s <- sweep(sweep(x_valid, 2, ctr, "-"), 2, scl, "/")
x_test_s  <- sweep(sweep(x_test,  2, ctr, "-"), 2, scl, "/")

mx <- suppressWarnings(max(abs(x_train_s)))
if (!is.finite(mx) || mx == 0) mx <- 1

x_train <- x_train_s / mx
x_valid <- x_valid_s / mx
x_test  <- x_test_s  / mx

# ------------------------------------------------------------
# 6) Run DDESONN
# ------------------------------------------------------------
res <- ddesonn_run(
  x = x_train,
  y = y_train,
  classification_mode = "binary",

  hidden_sizes = c(64, 32),
  seeds = 1L,
  do_ensemble = FALSE,

  validation = list(
    x = x_valid,
    y = y_valid
  ),

  test = list(
    x = x_test,
    y = y_test
  ),

  training_overrides = list(
    init_method = "he",
    optimizer = "adagrad",
    lr = 0.125,
    lambda = 0.00028,

    activation_functions = list(relu, relu, sigmoid),
    dropout_rates = list(0.10),
    loss_type = "CrossEntropy",

    validation_metrics = TRUE,
    num_epochs = 360,
    final_summary_decimals = 6L
  ),

  plot_controls = list(
    evaluate_report = list(
      roc_curve = TRUE,
      pr_curve  = FALSE
    )
  )
)
}

\end{ExampleCode}
\end{Examples}
\HeaderA{ddesonn\_training\_defaults}{Construct default training controls}{ddesonn.Rul.training.Rul.defaults}
%
\begin{Description}
Build a list of training hyperparameters that mirror the
expectations of the legacy DDESONN training loop.
\end{Description}
%
\begin{Usage}
\begin{verbatim}
ddesonn_training_defaults(
  mode = c("binary", "multiclass", "regression"),
  hidden_sizes = NULL
)
\end{verbatim}
\end{Usage}
%
\begin{Arguments}
\begin{ldescription}
\item[\code{mode}] Problem mode used to determine sensible defaults.

\item[\code{hidden\_sizes}] Integer vector describing the hidden layer widths.
\end{ldescription}
\end{Arguments}
%
\begin{Value}
A named list that can be modified and supplied to \code{\LinkA{ddesonn\_fit()}{ddesonn.Rul.fit}}.
\end{Value}
%
\begin{Examples}
\begin{ExampleCode}
ddesonn_training_defaults("binary", hidden_sizes = c(32, 16))

# Inspect regression defaults (includes LR decay by default).
cfg_reg <- ddesonn_training_defaults("regression", hidden_sizes = c(16, 8))
cfg_reg$lr
cfg_reg$lr_decay_rate
cfg_reg$lr_decay_epoch
cfg_reg$lr_min

# If you prefer a fixed LR in regression, disable decay explicitly.
cfg_reg$lr_decay_rate <- 1.0

\end{ExampleCode}
\end{Examples}
\HeaderA{predict.ddesonn\_model}{Predict method for DDESONN models}{predict.ddesonn.Rul.model}
%
\begin{Description}
This is the canonical user-facing wrapper around \code{\LinkA{ddesonn\_predict()}{ddesonn.Rul.predict}}. It
standardizes arguments (\code{type}, \code{aggregate}, \code{threshold}) and normalizes the
output structure for inference workflows.
Multiclass note: For multiclass classification, y should be encoded as integer class indices 1..K (or a one-hot matrix whose columns follow the model’s class order), otherwise accuracy comparisons may be incorrect.
\end{Description}
%
\begin{Usage}
\begin{verbatim}
## S3 method for class 'ddesonn_model'
predict(
  object,
  newdata,
  ...,
  aggregate = c("mean", "median", "none"),
  type = c("response", "class"),
  threshold = NULL,
  verbose = FALSE,
  verboseLow = FALSE,
  debug = FALSE
)
\end{verbatim}
\end{Usage}
%
\begin{Arguments}
\begin{ldescription}
\item[\code{object}] A ddesonn\_model (R6) returned by ddesonn\_run() / ddesonn\_model().

\item[\code{newdata}] Matrix/data.frame of predictors.

\item[\code{...}] Unused.

\item[\code{aggregate}] Aggregation strategy across ensemble members.

\item[\code{type}] Prediction type. \code{"response"} returns numeric predictions,
while \code{"class"} returns class labels for classification problems.

\item[\code{threshold}] Optional threshold override when \code{type = "class"}.

\item[\code{verbose}] Logical; emit detailed progress output when TRUE.

\item[\code{verboseLow}] Logical; emit important progress output when TRUE.

\item[\code{debug}] Logical; emit debug diagnostics when TRUE.
\end{ldescription}
\end{Arguments}
\HeaderA{print.ddesonn\_run\_result}{Print a summary of a DDESONN run result}{print.ddesonn.Rul.run.Rul.result}
%
\begin{Description}
Print a summary of a DDESONN run result
\end{Description}
%
\begin{Usage}
\begin{verbatim}
## S3 method for class 'ddesonn_run_result'
print(x, ...)
\end{verbatim}
\end{Usage}
%
\begin{Arguments}
\begin{ldescription}
\item[\code{x}] A \code{ddesonn\_run\_result} object.

\item[\code{...}] Unused.
\end{ldescription}
\end{Arguments}
%
\begin{Value}
\code{x}, invisibly.
\end{Value}
\printindex{}
\end{document}
